\section{Conclusion}
\label{sec:conclusion}

This paper has given a complete account of Burgers' equation, from its
origin as a simplified model of nonlinear advection and viscous diffusion
(Section~\ref{sec:motivation}) to the theory required to solve it once
classical solutions cease to exist. The method of characteristics
(Section~\ref{sec:characteristics}) gives an explicit solution formula
that is exact but necessarily temporary: Section~\ref{sec:shock-formation}
showed that decreasing initial data produces a genuine, finite-time loss
of smoothness, with an exact formula for the time and location of the
first singularity. Making sense of the equation beyond this point
required weakening the notion of solution
(Section~\ref{sec:weak-solutions}), which in turn revealed a genuine
failure of uniqueness (Section~\ref{sec:nonuniqueness}): the same initial
data can evolve into more than one weak solution. The Lax entropy
condition (Section~\ref{sec:entropy-condition}) was introduced precisely
to resolve this, and Section~\ref{sec:riemann-problem} used it to give
the complete, explicit entropy solution of the Riemann problem in both
possible cases.

The extensions in Section~\ref{sec:extensions} showed that this core
theory extends well beyond the specific flux $u^2/2$. The Cole--Hopf
transformation justified the entropy condition itself, as the selection
principle recovered in the vanishing-viscosity limit rather than an
arbitrary rule, and showed that viscosity alone, without the
vortex-stretching mechanism absent from this one-dimensional model, is
enough to guarantee global smoothness ---
supporting Remark~\ref{rem:navier-stokes-relevance}. The general convex
flux extension confirmed none of the core results are special to
Burgers' particular flux. The Godunov scheme, built directly on the
exact Riemann solver derived here, reproduced the analytic shock speed
and rarefaction profile numerically to within first-order error. The
shock interaction result showed multiple shocks merge via the same
Riemann-problem mechanism used throughout, the damped equation showed
this mechanism can be delayed or prevented by sufficiently strong linear
dissipation, and the N-wave result showed that the long-time behaviour of
any compactly supported solution is governed by a single universal
profile, determined only by the total positive and negative area of the
initial data.

Taken together, these results illustrate why Burgers' equation, despite
its simplicity, remains a standard reference point in the theory of
nonlinear partial differential equations: it is simple enough to solve
exactly and completely, in every regime considered here, while already
exhibiting the central phenomena --- shock formation, non-uniqueness of
weak solutions, the necessity of an entropy condition, and the
regularising effect of viscosity --- that make the general theory of
nonlinear conservation laws, and by extension the far harder open
questions surrounding the Navier--Stokes and Euler equations, so
difficult to resolve.